% ============================================================================
% MUSTERLOSUNG: Hacer planes - Plane machen
% Spanisch Klasse 10 | A1 | 90 Minuten
% ============================================================================

\documentclass[11pt,a4paper]{article}

\usepackage[utf8]{inputenc}
\usepackage[T1]{fontenc}
\usepackage[ngerman]{babel}
\usepackage[left=2cm,right=2cm,top=1.8cm,bottom=1.5cm]{geometry}
\usepackage{helvet}
\usepackage[table]{xcolor}
\usepackage{tabularx}
\usepackage{array}
\usepackage{enumitem}
\usepackage{tcolorbox}
\usepackage{fancyhdr}
\usepackage{lastpage}

\renewcommand{\familydefault}{\sfdefault}

\definecolor{spanisch}{HTML}{c41e3a}
\definecolor{grau}{HTML}{666666}
\definecolor{hellgrau}{HTML}{f5f5f5}
\definecolor{loesung}{HTML}{c41e3a}

\pagestyle{fancy}
\fancyhf{}
\fancyhead[L]{\small\textcolor{grau}{Spanisch Klasse 10 -- Hacer planes}}
\fancyhead[R]{\small\textcolor{loesung}{\textbf{MUSTERLOSUNG}}}
\fancyfoot[C]{\small\thepage/\pageref{LastPage}}
\renewcommand{\headrulewidth}{0.5pt}

\newcommand{\luecke}[1]{\underline{\hspace{#1}}}
\newcommand{\lsg}[1]{\textcolor{loesung}{\textbf{#1}}}
\newcommand{\punktebox}[1]{\hfill\fbox{\small \_\_/#1}}
\newcommand{\es}[1]{\textit{#1}}

\newtcolorbox{merkkasten}[1][]{
    colback=white,
    colframe=spanisch,
    fonttitle=\bfseries\color{spanisch},
    title=#1,
    boxrule=1.5pt,
    arc=0pt,
    left=5pt, right=5pt, top=5pt, bottom=5pt
}

\newtcolorbox{vokabelkasten}[1][]{
    colback=hellgrau,
    colframe=grau,
    fonttitle=\bfseries,
    title=#1,
    boxrule=0.5pt,
    arc=0pt,
    left=5pt, right=5pt, top=3pt, bottom=3pt
}

\newcounter{aufgabe}
\newcommand{\aufgabe}[2]{%
    \stepcounter{aufgabe}%
    \vspace{0.8em}%
    \noindent\textbf{\theaufgabe.} #1 \punktebox{#2}%
    \vspace{0.3em}%
}

\begin{document}

% ============================================================================
% KOPFBEREICH
% ============================================================================
\noindent
\begin{tabularx}{\textwidth}{@{}Xr@{}}
    {\Large\textbf{\textcolor{spanisch}{Hacer planes -- Plane machen}}} & \textcolor{loesung}{\Large\textbf{MUSTERLOSUNG}} \\[0.3em]
    {\small Spanisch Klasse 10 | A1 | 90 Minuten} & \\
\end{tabularx}

\vspace{0.3em}
\hrule
\vspace{0.6em}

% ============================================================================
% KASTEN V1: VOKABELN
% ============================================================================
\begin{vokabelkasten}[Kasten V1: Freizeitorte]
\begin{tabularx}{\textwidth}{@{}XXX@{}}
\es{el cine} -- das Kino & \es{la piscina} -- das Schwimmbad & \es{el polideportivo} -- das Sportzentrum \\
\es{el museo} -- das Museum & \es{el concierto} -- das Konzert & \es{el barrio} -- das Viertel \\
\end{tabularx}
\end{vokabelkasten}

\vspace{0.5em}

% ============================================================================
% KASTEN G1: IR - MIT LOSUNGEN
% ============================================================================
\begin{merkkasten}[Kasten G1: El verbo ir (gehen)]
\vspace{0.2em}
\textbf{Struktur:} ir + a + Ort \quad | \quad ir + al (m.) / a la (f.)

\vspace{0.3em}
\begin{tabularx}{\textwidth}{@{}|c|X|c|X|@{}}
\hline
\rowcolor{hellgrau}
\textbf{Person} & \textbf{ir} & \textbf{Person} & \textbf{ir} \\
\hline
yo & \lsg{voy} & nosotros/-as & \lsg{vamos} \\
\hline
tu & \lsg{vas} & vosotros/-as & \lsg{vais} \\
\hline
el/ella/usted & \lsg{va} & ellos/ellas/ustedes & \lsg{van} \\
\hline
\end{tabularx}
\end{merkkasten}

\vspace{0.3em}

% ============================================================================
% KASTEN G2: ADONDE
% ============================================================================
\begin{merkkasten}[Kasten G2: Die Frage ¿Adonde...?]
\es{¿Adonde vas?} = Wohin gehst du? \quad | \quad \es{¿Adonde vais el sabado?} = Wohin geht ihr am Samstag?
\end{merkkasten}

\vspace{0.3em}

% ============================================================================
% KASTEN G3: MODALVERBEN - MIT LOSUNGEN
% ============================================================================
\begin{merkkasten}[Kasten G3: Modalverben (querer / poder / tener que)]
\vspace{0.2em}
\textbf{Struktur:} Modalverb + Infinitiv \quad Beispiel: \es{Quiero ir al cine.}

\vspace{0.3em}
\begin{tabularx}{\textwidth}{@{}|c|X|X|X|@{}}
\hline
\rowcolor{hellgrau}
& \textbf{querer} (wollen) & \textbf{poder} (konnen) & \textbf{tener que} (mussen) \\
\hline
yo & \lsg{quiero} & \lsg{puedo} & tengo que \\
\hline
tu & quieres & \lsg{puedes} & tienes que \\
\hline
el/ella & \lsg{quiere} & puede & \lsg{tiene que} \\
\hline
nosotros & queremos & \lsg{podemos} & tenemos que \\
\hline
vosotros & \lsg{quereis} & podeis & \lsg{teneis que} \\
\hline
ellos & quieren & \lsg{pueden} & tienen que \\
\hline
\end{tabularx}
\end{merkkasten}

\vspace{0.3em}

% ============================================================================
% KASTEN G4: HAY
% ============================================================================
\begin{merkkasten}[Kasten G4: hay (es gibt)]
\es{hay} = es gibt (unveranderlich!)

\vspace{0.2em}
\begin{tabularx}{\textwidth}{@{}XX@{}}
\es{¿Hay un cine en tu barrio?} & Gibt es ein Kino in deinem Viertel? \\
\es{Si, hay un cine.} / \es{No, no hay cine.} & Ja, es gibt ein Kino. / Nein, es gibt kein Kino. \\
\end{tabularx}
\end{merkkasten}

\vspace{0.3em}

% ============================================================================
% KASTEN G5: GUSTAR
% ============================================================================
\begin{merkkasten}[Kasten G5: gustar (mogen/gefallen)]
\es{me gusta} = mir gefallt / ich mag \quad | \quad \es{te gusta} = dir gefallt / du magst

\vspace{0.2em}
\es{¿Te gusta ir al cine?} -- Gehst du gerne ins Kino? \quad | \quad \es{Si, me gusta.} -- Ja, ich mag das.
\end{merkkasten}

\vspace{0.8em}

% ============================================================================
% AUFGABE 1: UBERSETZUNG DE -> ES - MIT LOSUNGEN
% ============================================================================
\aufgabe{Ubersetze ins Spanische.}{10}

\begin{enumerate}[label=\alph*), leftmargin=2em, itemsep=0.6em]
    \item Ich gehe ins Kino.

    \lsg{Voy al cine.} \hfill \textcolor{grau}{(2 P.)}

    \item Willst du ins Schwimmbad gehen?

    \lsg{¿Quieres ir a la piscina?} \hfill \textcolor{grau}{(2 P.)}

    \item Wir konnen nicht gehen, weil wir lernen mussen.

    \lsg{No podemos ir porque tenemos que estudiar.} \hfill \textcolor{grau}{(2 P.)}

    \item Gibt es ein Sportzentrum in deinem Viertel?

    \lsg{¿Hay un polideportivo en tu barrio?} \hfill \textcolor{grau}{(2 P.)}

    \item Wohin gehst du am Samstag?

    \lsg{¿Adonde vas el sabado?} \hfill \textcolor{grau}{(2 P.)}
\end{enumerate}

\vspace{0.5em}

% ============================================================================
% AUFGABE 2: UBERSETZUNG ES -> DE - MIT LOSUNGEN
% ============================================================================
\aufgabe{Ubersetze ins Deutsche.}{10}

\begin{enumerate}[label=\alph*), leftmargin=2em, itemsep=0.6em]
    \item \es{¿Quieres ir al museo conmigo?}

    \lsg{Willst du mit mir ins Museum gehen?} \hfill \textcolor{grau}{(2 P.)}

    \item \es{No puedo ir porque tengo que trabajar.}

    \lsg{Ich kann nicht gehen, weil ich arbeiten muss.} \hfill \textcolor{grau}{(2 P.)}

    \item \es{¿Te gusta ir a conciertos?}

    \lsg{Gehst du gerne zu Konzerten? / Magst du es, zu Konzerten zu gehen?} \hfill \textcolor{grau}{(2 P.)}

    \item \es{Si, hay una piscina muy grande en mi barrio.}

    \lsg{Ja, es gibt ein sehr gro{\ss}es Schwimmbad in meinem Viertel.} \hfill \textcolor{grau}{(2 P.)}

    \item \es{Ellos van al polideportivo el domingo.}

    \lsg{Sie gehen am Sonntag ins Sportzentrum.} \hfill \textcolor{grau}{(2 P.)}
\end{enumerate}

\vspace{0.5em}

% ============================================================================
% AUFGABE 3: LUCKENTEXT DIALOG - MIT LOSUNGEN
% ============================================================================
\aufgabe{Erganze den Dialog mit den richtigen Verbformen.}{8}

\vspace{0.3em}
\noindent\textit{Ana trifft Carlos auf der Stra{\ss}e:}

\vspace{0.2em}
\begin{tabularx}{\textwidth}{@{}rX@{}}
\textbf{Ana:} & Hola Carlos, ¿\lsg{adonde vas} (adonde/ir, tu) este fin de semana? \\[0.5em]
\textbf{Carlos:} & \lsg{Voy} (ir, yo) al cine. ¿\lsg{Quieres} (querer, tu) venir conmigo? \\[0.5em]
\textbf{Ana:} & Si, \lsg{quiero} (querer, yo) ir. ¿A que hora? \\[0.5em]
\textbf{Carlos:} & ¿\lsg{Puedes} (poder, tu) ir a las cinco? \\[0.5em]
\textbf{Ana:} & No \lsg{puedo} (poder, yo) a las cinco porque \lsg{tengo que} (tener que, yo) estudiar. \\[0.5em]
\textbf{Carlos:} & Vale, entonces \lsg{vamos} (ir, nosotros) a las siete. ¡Perfecto! \\
\end{tabularx}

\vspace{0.8em}

% ============================================================================
% AUFGABE 4: EIGENER DIALOG - BEISPIELLOSUNG
% ============================================================================
\aufgabe{Schreibe deinen eigenen Dialog (mindestens 6 Zeilen).}{12}

\vspace{0.2em}
\noindent\textit{Situation: Du bist neu in der Stadt. Ein Freund/eine Freundin ladt dich ein.}\\
\textit{Verwende: querer, poder, tener que, hay, ir, ¿te gusta...?}

\vspace{0.3em}
\noindent\textcolor{loesung}{\textbf{Beispiellosung:}}

\vspace{0.2em}
\begin{tabularx}{\textwidth}{@{}|X|@{}}
\hline
\\[0.2em]
\textbf{A:} \lsg{Hola, ¿quieres ir al cine el sabado?} \\[0.5em]
\textbf{B:} \lsg{Si, quiero ir. ¿Hay un cine en tu barrio?} \\[0.5em]
\textbf{A:} \lsg{Si, hay un cine muy grande. ¿Te gusta ver peliculas?} \\[0.5em]
\textbf{B:} \lsg{Si, me gusta mucho. ¿A que hora vamos?} \\[0.5em]
\textbf{A:} \lsg{¿Puedes ir a las cuatro?} \\[0.5em]
\textbf{B:} \lsg{No puedo a las cuatro porque tengo que estudiar.} \\[0.5em]
\textbf{A:} \lsg{Vale, ¿podemos ir a las seis?} \\[0.5em]
\textbf{B:} \lsg{Si, a las seis puedo. ¡Perfecto!} \\[0.2em]
\hline
\end{tabularx}

\vspace{0.3em}
\noindent\textcolor{grau}{\textit{Bewertungskriterien: 6+ Zeilen (3 P.), alle Modalverben verwendet (3 P.), hay verwendet (2 P.), gustar verwendet (2 P.), sprachliche Korrektheit (2 P.)}}

\vspace{1em}
\hrule
\vspace{0.3em}
\begin{flushright}
\textbf{Gesamt: \luecke{1cm} / 40 Punkte}
\end{flushright}

\end{document}
